\documentclass[11pt]{article}
% ===========================================================================
%  Shared preamble for all MPM2D worksheets — input right after \documentclass.
%  (Files beginning with "_" are skipped by mpm2d-worksheet-publish.mjs.)
% ===========================================================================
\usepackage[letterpaper,top=2.2cm,bottom=1.9cm,left=1.6cm,right=1.6cm,headheight=28pt,footskip=24pt]{geometry}
\usepackage{amsmath,amssymb}
\usepackage[T1]{fontenc}
\usepackage[utf8]{inputenc}
\usepackage{lmodern}
\usepackage{xcolor}
\usepackage[most]{tcolorbox}
\usepackage{enumitem}
\usepackage{multicol}
\usepackage{fancyhdr}
\usepackage{tikz}
\usetikzlibrary{calc}
\usepackage{pgfplots}
\pgfplotsset{compat=1.18}

\definecolor{exblue}{HTML}{4A90E2}
\definecolor{exbg}{HTML}{E6F3FF}
\definecolor{qorange}{HTML}{E69138}
\definecolor{qbg}{HTML}{FFF7CC}
\definecolor{keygray}{HTML}{475569}
\definecolor{keybg}{HTML}{F1F5F9}
\definecolor{iagreen}{HTML}{14653B}
\definecolor{titleblue}{HTML}{1D4ED8}
% Rotating palette so each worked example gets its own colour.
\definecolor{ex0}{HTML}{2563A0}\definecolor{ex1}{HTML}{3B7D3B}\definecolor{ex2}{HTML}{8A5A00}
\definecolor{ex3}{HTML}{A3327A}\definecolor{ex4}{HTML}{6D28A3}\definecolor{ex5}{HTML}{B08900}
\definecolor{ex6}{HTML}{0E7490}\definecolor{ex7}{HTML}{9A3412}\definecolor{ex8}{HTML}{15803D}

\pagestyle{fancy}
\fancyhf{}
\fancyhead[L]{\footnotesize NAME: \underline{\hspace{3.4cm}}}
\fancyhead[C]{\textbf{Integration Academy}\\[-2pt]{\scriptsize Math Simplified \textperiodcentered{} Success Amplified}}
\fancyhead[R]{\footnotesize MPM2D}
\fancyfoot[L]{\scriptsize dr.merie@IntegrationAcademy.ca}
\fancyfoot[C]{\scriptsize\thepage}
\fancyfoot[R]{\scriptsize IntegrationAcademy.ca}
\renewcommand{\headrulewidth}{1.2pt}
\renewcommand{\footrulewidth}{0.4pt}
\setlength{\parindent}{0pt}
\setlength{\parskip}{4pt}

% Examples are NOT breakable, so each example + its solution stays on one page.
\newtcolorbox{example}[2]{colframe=#1,colback=#1!7,coltitle=white,colbacktitle=#1,fonttitle=\bfseries,title={#2},arc=3pt,boxrule=1pt}
\newtcolorbox{qbox}[1]{colframe=qorange,colback=qbg,coltitle=white,colbacktitle=qorange,fonttitle=\bfseries,title={#1},arc=3pt,breakable,boxrule=1pt}
\newtcolorbox{keybox}{colframe=keygray,colback=keybg,coltitle=white,colbacktitle=keygray,fonttitle=\bfseries,title={$\checkmark$\ Answer Key},arc=3pt,breakable,boxrule=1pt}
\newtcolorbox{recapbox}{colframe=keygray,colback=keybg,coltitle=white,colbacktitle=keygray,fonttitle=\bfseries,title={Question Recap},arc=3pt,breakable,boxrule=1pt}
\newtcolorbox{ideabox}{colframe=titleblue,colback=titleblue!6,coltitle=white,colbacktitle=titleblue,fonttitle=\bfseries,title={The Big Ideas},arc=3pt,breakable,boxrule=1.2pt}
\newtcolorbox{learnbox}{colframe=iagreen,colback=iagreen!5,coltitle=white,colbacktitle=iagreen,fonttitle=\bfseries,title={Concept Essentials},arc=3pt,breakable,boxrule=1.2pt}
\newcommand{\lh}[1]{\par\smallskip{\bfseries\color{iagreen}#1}\par\nobreak\smallskip}

\newcommand{\soln}{\par\smallskip\textit{\textbf{Solution.}}\ }
\newcommand{\workspace}[1]{\par\vspace{3pt}\fcolorbox{black!30}{white}{\begin{minipage}[t][#1][t]{\dimexpr\linewidth-2\fboxsep\relax}\ \end{minipage}}\par}
\newcommand{\sech}[1]{\par\vspace{8pt}{\Large\bfseries\color{iagreen}#1}\par\nobreak\vspace{2pt}{\color{black!20}\hrule height1pt}\vspace{6pt}}

% A compact coordinate plot sized for an example box: \eplot{xmin}{xmax}{ymin}{ymax}{body}
\newcommand{\eplot}[5]{\begin{center}\begin{tikzpicture}
\begin{axis}[width=6.8cm,height=4.4cm,axis lines=middle,xlabel={\tiny$x$},ylabel={\tiny$y$},
grid=both,grid style={gray!16},xmin=#1,xmax=#2,ymin=#3,ymax=#4,
xtick distance=2,ytick distance=2,tick label style={font=\tiny},axis line style={-},samples=120]
#5
\end{axis}\end{tikzpicture}\end{center}}

% A sine-curve plot (degrees on x, 0..360): \splot{ymin}{ymax}{body}
\newcommand{\splot}[3]{\begin{center}\begin{tikzpicture}
\begin{axis}[width=8.6cm,height=4.6cm,axis lines=middle,xlabel={\tiny$x\,(^\circ)$},ylabel={\tiny$y$},
grid=both,grid style={gray!16},xmin=0,xmax=360,ymin=#1,ymax=#2,
xtick distance=90,ytick distance=1,tick label style={font=\tiny},axis line style={-},samples=180]
#3
\end{axis}\end{tikzpicture}\end{center}}

% A small coordinate plot: \plot{xmin}{xmax}{ymin}{ymax}{body}
\newcommand{\plot}[5]{\begin{center}\begin{tikzpicture}
\begin{axis}[width=8.4cm,height=6cm,axis lines=middle,xlabel={\scriptsize$x$},ylabel={\scriptsize$y$},
grid=both,grid style={gray!16},xmin=#1,xmax=#2,ymin=#3,ymax=#4,
xtick distance=2,ytick distance=2,tick label style={font=\tiny},axis line style={-}]
#5
\end{axis}\end{tikzpicture}\end{center}}

% Right triangle (angle theta at bottom-left, right angle at bottom-right):
%   \rtri{drawBase}{drawHeight}{oppLabel}{adjLabel}{hypLabel}
\newcommand{\rtri}[5]{\begin{center}\begin{tikzpicture}[scale=0.85]
\coordinate (A) at (0,0);\coordinate (B) at (#1,0);\coordinate (C) at (#1,#2);
\draw[thick,exblue] (A)--(B)--(C)--cycle;
\draw[black] ($(B)+(-0.32,0)$)--($(B)+(-0.32,0.32)$)--($(B)+(0,0.32)$);
\node[below] at ($(A)!0.5!(B)$){#4};
\node[right] at ($(B)!0.5!(C)$){#3};
\node[above left] at ($(A)!0.5!(C)$){#5};
\node at (0.7,0.22){$\theta$};
\end{tikzpicture}\end{center}}

% General (oblique) triangle with vertex labels and opposite-side labels:
%   \gtri{Alabel}{Blabel}{Clabel}{aSide}{bSide}{cSide}
\newcommand{\gtri}[6]{\begin{center}\begin{tikzpicture}[scale=0.85]
\coordinate (A) at (0,0);\coordinate (B) at (5,0);\coordinate (C) at (1.6,3);
\draw[thick,exblue] (A)--(B)--(C)--cycle;
\node[below left] at (A){#1};\node[below right] at (B){#2};\node[above] at (C){#3};
\node[below] at ($(A)!0.5!(B)$){#6};
\node[right] at ($(B)!0.5!(C)$){#4};
\node[left] at ($(A)!0.5!(C)$){#5};
\end{tikzpicture}\end{center}}

\fancyhead[R]{\footnotesize MCR3U \textperiodcentered{} 1.4}
\begin{document}

\begin{center}
{\LARGE\bfseries\color{titleblue}Worksheet 1.4 --- Transformations of Functions}\\[2pt]
{\small Grade 11 Functions, University (MCR3U) \textperiodcentered{} Unit 1: Characteristics of Functions}
\end{center}

\begin{ideabox}
Every transformed function has the form $g(x)=a\,f\big(k(x-d)\big)+c$.
\begin{itemize}[leftmargin=*,itemsep=2pt]
  \item $a$: vertical stretch/reflection. $c$: vertical shift.
  \item $d$: horizontal shift (opposite to the sign inside).
  \item Apply stretches/reflections first, then translations.
\end{itemize}
\end{ideabox}

\begin{learnbox}
\lh{The general form}In $y=af\!\left(k(x-d)\right)+c$, each parameter moves the base graph: $a$ stretches vertically, $k$ stretches horizontally, $d$ shifts right, $c$ shifts up.
\lh{Order and signs}Mind the minus: $(x-d)$ shifts \emph{right} by $d$. A negative $a$ reflects in the $x$-axis; a negative $k$ reflects in the $y$-axis.
\lh{Mapping points}Apply the changes to key points of the base function: $(x,y)\to\left(\dfrac{x}{k}+d,\ ay+c\right)$ places each on the new graph.
\end{learnbox}

\sech{Worked Examples}

\begin{example}{ex0}{Example 1 --- Full description}
Describe $y=-2(x-3)^2+1$ from $y=x^2$.\soln $a=-2$: reflect and stretch by 2. $d=3$: right 3. $c=1$: up 1. Vertex $(3,1)$, opens down:\eplot{-1}{7}{-8}{3}{\addplot[exblue,very thick,domain=0.7:5.3]{-2*(x-3)^2+1};\addplot[red,only marks,mark=*,mark size=1.6pt] coordinates {(3,1)};}
\end{example}

\begin{example}{ex1}{Example 2 --- Radical}
Describe $y=\sqrt{x+4}-2$.\soln $d=-4$: left 4. $c=-2$: down 2.
\end{example}

\begin{example}{ex2}{Example 3 --- Reflection}
Describe $y=-x^2$.\soln $a=-1$: a reflection in the $x$-axis (opens down).
\end{example}

\begin{example}{ex3}{Example 4 --- Stretch}
Describe $y=3x^2$.\soln $a=3$: a vertical stretch by 3.
\end{example}

\begin{example}{ex4}{Example 5 --- Vertex from form}
Vertex of $y=(x+1)^2-5$?\soln $h=-1$, $k=-5$: vertex $(-1,-5)$.
\end{example}

\begin{example}{ex5}{Example 6 --- Description}
Describe $y=(x-2)^2+3$.\soln Right 2, up 3.
\end{example}

\begin{example}{ex6}{Example 7 --- Radical}
Describe $y=\sqrt{x}-4$.\soln Down 4.
\end{example}

\begin{example}{ex7}{Example 8 --- Vertex}
Vertex of $y=-(x-4)^2+2$?\soln $(4,2)$; $a<0$ so it opens down.
\end{example}

\begin{example}{ex8}{Example 9 --- Stretch + shift}
Describe $y=2(x+1)^2$.\soln Stretch by 2, left 1.
\end{example}

\clearpage
\sech{Your Turn --- Show Your Work}
For each question, show all steps clearly in the space provided.

\begin{qbox}{Question 1}
Describe $y=(x-2)^2+3$.\workspace{2.4cm}
\end{qbox}

\begin{qbox}{Question 2}
Describe $y=\sqrt{x}-4$.\workspace{2.4cm}
\end{qbox}

\begin{qbox}{Question 3}
State the vertex of $y=-(x-4)^2+2$.\workspace{2.4cm}
\end{qbox}

\begin{qbox}{Question 4}
What does $a=-1$ do?\workspace{2.4cm}
\end{qbox}

\begin{qbox}{Question 5}
Describe $y=2(x+1)^2$.\workspace{2.4cm}
\end{qbox}

\begin{qbox}{Question 6}
State the vertex of $y=(x+1)^2-5$.\workspace{2.4cm}
\end{qbox}

\begin{qbox}{Question 7}
Describe $y=-2(x-3)^2+1$.\workspace{2.4cm}
\end{qbox}

\begin{qbox}{Question 8}
Write $y=x^2$ shifted right $4$ and down $2$.\workspace{2.4cm}
\end{qbox}

\begin{qbox}{Question 9}
Describe $y=\sqrt{x+4}-2$.\workspace{2.4cm}
\end{qbox}

\begin{qbox}{Question 10}
Does $y=-3(x+2)^2+4$ open up or down?\workspace{2.4cm}
\end{qbox}

\begin{qbox}{Question 11}
State the vertex of $y=(x-5)^2+2$.\workspace{2.4cm}
\end{qbox}

\begin{qbox}{Question 12}
What does $(x-3)$ inside do?\workspace{2.4cm}
\end{qbox}

\begin{qbox}{Question 13 --- Challenge}
Describe $y=-(x+1)^2+5$ fully.\workspace{3cm}
\end{qbox}

\clearpage
\sech{Question Recap \& Answer Key}

\begin{recapbox}
\begin{multicols}{2}
\small
\begin{enumerate}[leftmargin=*,itemsep=3pt]
  \item Describe $y=(x-2)^2+3$.
  \item Describe $y=\sqrt{x}-4$.
  \item State the vertex of $y=-(x-4)^2+2$.
  \item What does $a=-1$ do?
  \item Describe $y=2(x+1)^2$.
  \item State the vertex of $y=(x+1)^2-5$.
  \item Describe $y=-2(x-3)^2+1$.
  \item Write $y=x^2$ shifted right $4$ and down $2$.
  \item Describe $y=\sqrt{x+4}-2$.
  \item Does $y=-3(x+2)^2+4$ open up or down?
  \item State the vertex of $y=(x-5)^2+2$.
  \item What does $(x-3)$ inside do?
  \item Describe $y=-(x+1)^2+5$ fully.
\end{enumerate}
\end{multicols}
\end{recapbox}

\begin{keybox}
\begin{multicols}{2}
\small
\begin{enumerate}[leftmargin=*,itemsep=3pt]
  \item right $2$, up $3$
  \item down $4$
  \item $(4,2)$
  \item reflects in the $x$-axis
  \item stretch $2$, left $1$
  \item $(-1,-5)$
  \item stretch $2$, reflect, right $3$, up $1$
  \item $y=(x-4)^2-2$
  \item left $4$, down $2$
  \item down
  \item $(5,2)$
  \item shifts right $3$
  \item reflect in $x$-axis, left $1$, up $5$; vertex $(-1,5)$
\end{enumerate}
\end{multicols}
\end{keybox}

\end{document}
